\section{Conclusão}\label{sec:conclusao}

Decomposto em subsistemas, o modelo confirma a ordem de grandeza que a literatura atribui ao motor inteiro.
O pistão falha 62 a 83 vezes mais que os motores a turbina em um voo de 1~h, e a diferença resiste a toda a faixa de taxas publicadas.
Nele, 60\% das falhas estão no núcleo mecânico, que não tem redundância, e é ali que manutenção e projeto têm mais efeito; a ignição duplicada protege uma função que já respondia por pouco das falhas.
Nos motores a turbina, compressores e turbinas concentram as falhas, e o FADEC de dois canais quase não aparece nos registros.
A redundância não compensa a tecnologia: dois motores a pistão só superam uma turbina se a causa comum entre eles ficar abaixo de 1,2\% a 1,6\%, e, quando a falha de um motor de um multimotor a pistão chega ao NTSB, 43\% dos casos têm vítima fatal.
Entre turboélice e turbofan, as faixas da literatura não permitem estabelecer uma ordem.
O teste caso-controle com o tempo desde a revisão sustenta a taxa constante no pistão.

Três frentes ficam em aberto.
Obter horas voadas por frota permitiria a atualização Gama--Poisson das taxas absolutas, que a falta de exposição impediu.
Ampliar a base das turbinas, que têm 69 e 45 registros, estreitaria os intervalos do Pareto dessas tecnologias.
E taxas de remoção em campo das unidades redundantes, como magnetos e canais do FADEC, dariam um limite superior para $\lambda_u$, já que o NTSB só vê as falhas que terminam em acidente.
